\textbf{结论名称：} 两相交圆公共弦所在直线的方程\\
\textbf{适用条件：} 两个不同圆 $C_1,C_2$ 交于两个不同点 $A,B$，圆心为 $O_1(a_1,b_1)$、$O_2(a_2,b_2)$，半径为 $r_1,r_2>0$，圆心距 $d>0$，且 $|r_1-r_2|<d<r_1+r_2$。\\
\textbf{变量说明：} 令 $M(x_M,y_M)$ 为公共弦 $AB$ 的中点；$h=AM=BM$；令 $t$ 为沿 $O_1\to O_2$ 方向从 $O_1$ 到 $M$ 的有向距离，$\Delta a=a_2-a_1$、$\Delta b=b_2-b_1$，$d=\sqrt{(\Delta a)^2+(\Delta b)^2}$。\\
\textbf{核心公式：}
\[
\boxed{t=\frac{d^2+r_1^2-r_2^2}{2d}},\qquad
\boxed{M=(a_1,b_1)+\frac{t}{d}(\Delta a,\Delta b)}
\]
\[
\boxed{\Delta a(x-x_M)+\Delta b(y-y_M)=0}
\]
等价地，$\boxed{2\Delta a(x-a_1)+2\Delta b(y-b_1)=d^2+r_1^2-r_2^2}$。若 $C_i:x^2+y^2+D_ix+E_iy+F_i=0$，则直接相减得 $\boxed{(D_1-D_2)x+(E_1-E_2)y+F_1-F_2=0}$。\\
\textbf{使用场景：} 求两相交圆公共弦所在直线、弦中点、弦长以及解释弦的位置；只求直线可直接相减，求位置与长度可优先采用几何法。\\
\textbf{简单例子：} $C_1:x^2+y^2=25$，$C_2:(x-3)^2+(y-4)^2=16$。连心线斜率 $4/3$，公共弦斜率 $-3/4$；由 $t=(25+25-16)/10=17/5$ 得 $M=(51/25,68/25)$，于是 $y-68/25=-\frac34(x-51/25)$，即 $\boxed{3x+4y-17=0}$。两圆方程相减也得到 $6x+8y-34=0$。
